\documentclass[11pt]{article}

\usepackage[T1]{fontenc}
\usepackage[utf8]{inputenc}
\usepackage{lmodern}
\usepackage{microtype}
\usepackage[margin=1in]{geometry}
\usepackage{amsmath,amssymb,amsthm,mathtools,bm}
\usepackage{booktabs,longtable,array}
\usepackage{enumitem}
\usepackage{xcolor}
\usepackage{hyperref}
\usepackage[nameinlink,capitalise,noabbrev]{cleveref}
\usepackage{fancyhdr}

\hypersetup{
  colorlinks=true,
  linkcolor=blue!55!black,
  citecolor=green!35!black,
  urlcolor=blue!60!black,
  pdftitle={Primitive exponent lattices and Hankel orbits in the Alaoglu--Erdos problem},
  pdfauthor={Research progress report}
}

\pagestyle{fancy}
\fancyhf{}
\fancyhead[L]{Alaoglu--Erd\H{o}s progress report}
\fancyhead[R]{Primitive lattices and Hankel orbits}
\fancyfoot[C]{\thepage}
\setlength{\headheight}{14pt}

\numberwithin{equation}{section}
\setlist[itemize]{topsep=4pt,itemsep=2pt,parsep=1pt}
\setlist[enumerate]{topsep=4pt,itemsep=2pt,parsep=1pt}

\newtheorem{theorem}{Theorem}[section]
\newtheorem{proposition}[theorem]{Proposition}
\newtheorem{lemma}[theorem]{Lemma}
\newtheorem{corollary}[theorem]{Corollary}
\newtheorem{criterion}[theorem]{Criterion}
\newtheorem{conjecture}[theorem]{Conjecture}
\theoremstyle{definition}
\newtheorem{definition}[theorem]{Definition}
\newtheorem{example}[theorem]{Example}
\newtheorem{remark}[theorem]{Remark}
\newtheorem{warning}[theorem]{Warning}

\newcommand{\Q}{\mathbb{Q}}
\newcommand{\Z}{\mathbb{Z}}
\newcommand{\N}{\mathbb{N}}
\newcommand{\R}{\mathbb{R}}
\newcommand{\C}{\mathbb{C}}
\newcommand{\Qbar}{\overline{\mathbb{Q}}}
\newcommand{\Gm}{\mathbb{G}_{\mathrm m}}
\newcommand{\vpp}{v_p}
\newcommand{\val}{\boldsymbol{v}}
\newcommand{\Span}{\operatorname{Span}}
\newcommand{\rank}{\operatorname{rank}}
\newcommand{\sat}{\operatorname{sat}}
\newcommand{\cont}{\operatorname{cont}}
\newcommand{\SNF}{\operatorname{SNF}}
\newcommand{\supp}{\operatorname{supp}}
\newcommand{\dd}{\,\mathrm d}
\newcommand{\Pair}{\mathbf P}
\newcommand{\axis}{\mathbf e}
\newcommand{\defect}{\delta}
\newcommand{\Ealg}{\mathcal E_{\mathrm{alg}}}
\newcommand{\Erat}{\mathcal E_{\mathrm{rat}}}
\newcommand{\Eint}{\mathcal E_{\mathrm{int}}}
\newcommand{\Msol}{\mathcal M}

\title{\textbf{Beyond the Two-by-Two Barrier}\\[4pt]
Primitive Exponent Lattices, Exact Descent, and Rational Hankel Orbits\\
for the Alaoglu--Erd\H{o}s Conjecture}
\author{A nonduplicative research progress report}
\date{23 August 2026}

\begin{document}
\maketitle

\begin{center}
\fbox{\begin{minipage}{0.92\textwidth}
\textbf{Status.} This report does \emph{not} claim a proof of the Alaoglu--Erd\H{o}s conjecture. It proves a collection of exact conditional structure theorems. The strongest one says that, if the conjecture is false, then after a canonical and finite root-extraction procedure there is a unique transcendental exponent \(\beta\) for which the entire set of common integral exponents is exactly the free monoid
\[
        \N_0+\N_0\beta.
\]
Every other nonintegral solution is an integer shift of a positive integral power of this one primitive orbit. A second exact result classifies every rational three-term Hankel orbit and shows that extending a putative counterexample by only one further exponentiation step would force rationality. The remaining gap is isolated in several equivalent, sharply stated forms.
\end{minipage}}
\end{center}

\begin{abstract}
Write
\[
  \Eint(2,3)=\{t\in\R:2^t,3^t\in\N\}.
\]
The Alaoglu--Erd\H{o}s conjecture asserts that \(\Eint(2,3)=\N_0\). The classical Four Exponentials conjecture would settle this immediately, whereas the unconditional Six Exponentials theorem leaves a one-dimensional gap. The aim here is to exploit that gap rather than merely restate it.

Assume, conditionally, that a nonintegral element exists. We introduce a finite valuation lattice generated by the axis pair \((2,3)\) and a counterexample pair \((A,B)=(2^x,3^x)\). Its saturation index is the explicit integer
\[
 \defect(A,B)=\gcd\!\left(
 \{v_p(A):p\ne2\},\{v_p(B):p\ne3\},v_3(B)-v_2(A)
 \right).
\]
This index is exactly the maximal common root that can be extracted after a matched axis shift. Dividing it out and removing the largest common axis factor produces a canonical pair \((A_0,B_0)=(2^\beta,3^\beta)\) satisfying \(\defect(A_0,B_0)=1\) and \(\min(v_2(A_0),v_3(B_0))=0\). Combining this primitive valuation normal form with Six Exponentials gives
\[
 \Erat(2,3)=\Z+\Z\beta,
 \qquad
 \Eint(2,3)=\N_0+\N_0\beta.
\]
The canonical \(\beta\) is unique and is the least nonintegral common integral exponent. If \(\alpha=n+k\beta\), then its axis content is \(n\), its saturation defect is exactly \(k\), and its localized pair is \((A_0^k,B_0^k)\). The associated determinant lift has exact Smith normal form \(\operatorname{diag}(1,k,k)\).

A separate orbit theorem combines Six Exponentials with Gelfond--Schneider: for multiplicatively independent algebraic bases \(a,b\), simultaneous algebraicity of
\[
 a,a^\xi,a^{\xi^2}
 \quad\text{and}\quad
 b,b^\xi,b^{\xi^2}
\]
forces \(\xi\in\Q\). In the rational case \(\xi=r/s\), every prime valuation of the three-term orbit has the exact pattern
\[
       d_p(s^2,rs,r^2),
\]
and the orbit obeys two monomial recurrences. This yields a complete classification of the localized rank-three Hankel branch, a reciprocal duality for the base pair \((A_0,B_0)\), and a list of concrete criteria whose proof would finish the conjecture. The report concludes with a formalization plan separating elementary lattice lemmas from the two transcendence inputs.
\end{abstract}

\tableofcontents
\newpage

\section{Scope, novelty, and the exact target}

The problem is deceptively short:
\begin{conjecture}[Alaoglu--Erd\H{o}s]
If \(x\in\R\) and \(2^x,3^x\in\Z\), then \(x\in\Z\).
\end{conjecture}
Because positive real powers are positive, the hypothesis really places both values in \(\N\). The purpose of this report is not to repeat the standard equivalence with a special case of Four Exponentials, nor to repeat broad determinant, interpolation, finite-search, or generic \(p\)-adic programs. Instead it develops four tightly connected points that are useful even without a final proof:

\begin{enumerate}
\item an exact valuation-lattice invariant measuring all possible rational root extraction;
\item a canonical primitive counterexample and a complete conditional description of every other common rational or integral exponent;
\item an exact rank-three orbit/Hankel theorem, including the square valuation pattern and its converse;
\item a Smith-normal-form and fractional-linear formalism that identifies which operations merely recycle the same two-dimensional exponent lattice and which operations would genuinely escape it.
\end{enumerate}

The central distinction throughout is between a \emph{proved reduction} and a \emph{desired closure property}. For example, proving that \(2^{x^2}\) and \(3^{x^2}\) are algebraic would finish the problem, but their algebraicity does not follow from the original hypotheses. The report records this as a precise one-step orbit barrier rather than silently assuming it.

\section{Elementary reductions and notation}

For positive numbers \(u,v\), write
\[
  \Pair_{u,v}(t)=(u^t,v^t).
\]
For the distinguished bases we abbreviate
\[
  \Pair(t)=\Pair_{2,3}(t)=(2^t,3^t),
  \qquad \axis=\Pair(1)=(2,3).
\]
We use three exponent sets:
\begin{align*}
 \Ealg(u,v)&=\{t\in\R:u^t,v^t\in\Qbar\},\\
 \Erat(u,v)&=\{t\in\R:u^t,v^t\in\Q_{>0}\},\\
 \Eint(u,v)&=\{t\in\R:u^t,v^t\in\N\}.
\end{align*}
Thus the conjecture is \(\Eint(2,3)=\N_0\).

\begin{lemma}[Sign and rational-exponent reductions]\label{lem:elementary}
Let \(x\in\Eint(2,3)\).
\begin{enumerate}[label=\textup{(\roman*)}]
\item \(x\ge0\).
\item If \(x\in\Q\), then \(x\in\Z\).
\item If \(x\notin\Z\), then \(x\ge1\), \(x\) is irrational, and in fact \(x\) is transcendental.
\end{enumerate}
\end{lemma}

\begin{proof}
If \(x<0\), then \(0<2^x<1\), so \(2^x\notin\N\). If \(x=a/b\) in lowest terms with \(b>0\), then \(2^a=(2^x)^b\). Unique factorization forces \(b\mid a\), hence \(b=1\). A positive nonzero integral value of \(2^x\) is at least \(2\), so \(x\ge1\). Finally, if a nonintegral \(x\) were algebraic, it would be algebraic irrational; Gelfond--Schneider applied to \(2^x\) would make \(2^x\) transcendental, contrary to the hypothesis.
\end{proof}

Fix for the moment a hypothetical nonintegral solution \(x\), and put
\[
  A=2^x,\qquad B=3^x.
\]
Then \(A,B\in\N\), and
\[
   \frac{\log A}{\log 2}=\frac{\log B}{\log 3}=x.
\]
The numbers \(A\) and \(B\) are multiplicatively independent. Indeed, \(A^r=B^s\) with nonzero integers \(r,s\) would imply
\(
 r\log A=s\log B
\), hence \(r\log2=s\log3\), contradicting unique factorization of \(2^r=3^s\).

For a positive rational \(q\), let
\[
  \val(q)=(v_p(q))_{p\ \mathrm{prime}}\in\Z^{(\mathcal P)},
\]
the finitely supported prime-valuation vector. For a rational pair write
\[
  \val(u,v)=\bigl(\val(u),\val(v)\bigr)
       \in \Z^{(\mathcal P)}\oplus\Z^{(\mathcal P)}.
\]
Set
\[
  g_0=\val(2,3),\qquad g_1=\val(A,B).
\]
The two vectors are \(\Q\)-linearly independent because \(1,x\) are.

\section{The Six Exponentials rank ceiling}

We use the following classical theorem in its standard real-positive specialization.

\begin{theorem}[Six Exponentials]\label{thm:sixexp}
Let \(z_1,z_2\in\C\) be linearly independent over \(\Q\), and let \(w_1,w_2,w_3\in\C\) be linearly independent over \(\Q\). Then at least one of the six numbers
\[
       \exp(z_iw_j),\qquad 1\le i\le2,\quad 1\le j\le3,
\]
is transcendental.
\end{theorem}

The first structural consequence is elementary but decisive.

\begin{proposition}[Common algebraic exponents have rank at most two]\label{prop:ranktwo}
If \(u,v>0\) are multiplicatively independent algebraic numbers, then \(\Ealg(u,v)\) is a \(\Q\)-vector space of dimension at most two.
\end{proposition}

\begin{proof}
Closure under addition and subtraction follows by multiplying and dividing algebraic values. If \(t\in\Ealg(u,v)\) and \(q\in\Q\), positive real choices of the relevant roots show that \(u^{qt}\) and \(v^{qt}\) are algebraic. Thus \(\Ealg(u,v)\) is a \(\Q\)-vector space.

If it contained three independent elements \(t_1,t_2,t_3\), apply \cref{thm:sixexp} with \(z_1=\log u\), \(z_2=\log v\), and \(w_j=t_j\). Multiplicative independence makes the two logarithms \(\Q\)-independent, while all six exponentials are algebraic by definition, a contradiction.
\end{proof}

\begin{corollary}[Exact algebraic exponent plane under a counterexample]\label{cor:Ealgplane}
If \(x\notin\Z\) and \(A=2^x,B=3^x\in\N\), then
\[
       \Ealg(2,3)=\Q+\Q x.
\]
\end{corollary}

\begin{proof}
The left-hand side contains the independent elements \(1,x\), so \cref{prop:ranktwo} gives containment in their \(\Q\)-span. Conversely, if \(r,s\in\Q\), then
\[
  2^{r+sx}=2^rA^s,\qquad 3^{r+sx}=3^rB^s
\]
are algebraic.
\end{proof}

This is the exact content supplied by Six Exponentials. It does not yet say which points of the plane give \emph{rational} values. That arithmetic distinction is where the valuation lattice enters.

\section{The valuation lattice and its saturation defect}

Let
\[
   L(A,B)=\Z g_0+\Z g_1
   \subset \Z^{(\mathcal P)}\oplus\Z^{(\mathcal P)}.
\]
Its saturation is
\[
  L(A,B)^{\sat}
  =\bigl(\Q L(A,B)\bigr)\cap
       \bigl(\Z^{(\mathcal P)}\oplus\Z^{(\mathcal P)}\bigr).
\]
Only finitely many coordinates occur, so all Smith-normal-form statements below concern an ordinary finite integer matrix obtained by deleting zero rows.

Write
\[
 a_p=v_p(A),\qquad b_p=v_p(B).
\]

\begin{definition}[Saturation defect]\label{def:defect}
Define
\begin{equation}\label{eq:defect}
 \defect(A,B)=\gcd\!\left(
   \{a_p:p\ne2\},
   \{b_p:p\ne3\},
   b_3-a_2
 \right),
\end{equation}
where zeros may be omitted and the gcd is taken positive.
\end{definition}

For a nonintegral solution the displayed collection contains a nonzero integer: otherwise \(A\) would be a power of \(2\) and \(B\) a power of \(3\), forcing \(x\in\Z\).

\begin{theorem}[Exact saturation index]\label{thm:satindex}
For a hypothetical counterexample pair \((A,B)\),
\[
   [L(A,B)^{\sat}:L(A,B)]=\defect(A,B).
\]
Equivalently, the two nonzero invariant factors of the valuation matrix with columns \(g_0,g_1\) are
\[
       1,\ \defect(A,B).
\]
\end{theorem}

\begin{proof}
For a rank-two integer matrix, the first determinantal divisor is the gcd of all entries and the second is the gcd of all \(2\times2\) minors. The first is \(1\), because \(g_0\) contains entries equal to \(1\).

The only potentially nonzero \(2\times2\) minors are obtained by pairing one of the two nonzero rows of \(g_0\) with another row. Pairing the first-coordinate row at \(p=2\) with a first-coordinate row at \(p\ne2\) gives \(a_p\); pairing it with the second-coordinate row at \(p=3\) gives \(b_3-a_2\); pairing it with a second-coordinate row at \(p\ne3\) gives \(b_p\). Pairing the row at the second-coordinate prime \(3\) with other rows only repeats these numbers up to sign. Their gcd is exactly \eqref{eq:defect}. The Smith normal form is therefore \(\operatorname{diag}(1,\defect(A,B))\), and the product of its nonunit invariant factors is the saturation index.
\end{proof}

The invariant has a direct root-extraction interpretation.

\begin{theorem}[Matched-axis root extraction]\label{thm:rootextract}
Let \(d\ge1\). The following are equivalent.
\begin{enumerate}[label=\textup{(\roman*)}]
\item \(d\mid\defect(A,B)\).
\item There is an integer \(r\), unique modulo \(d\), and positive integers \(C,D\) such that
\begin{equation}\label{eq:rootfactor}
       A=2^r C^d,\qquad B=3^r D^d.
\end{equation}
\item The classes of \(A\) and \(B\) in the Kummer quotients satisfy
\[
 [A]=[2]^r\ \text{in }\Q_{>0}^{\times}/\Q_{>0}^{\times d},
 \qquad
 [B]=[3]^r\ \text{in }\Q_{>0}^{\times}/\Q_{>0}^{\times d}
\]
for the same residue \(r\bmod d\).
\end{enumerate}
If these conditions hold and \(A=2^x,B=3^x\), then
\begin{equation}\label{eq:descentexp}
       y=\frac{x-r}{d}
\end{equation}
is again a common integral exponent, with \(2^y=C\) and \(3^y=D\).
\end{theorem}

\begin{proof}
Condition \(d\mid\defect(A,B)\) says that every \(a_p\) for \(p\ne2\) and every \(b_p\) for \(p\ne3\) is divisible by \(d\), and that \(a_2\equiv b_3\pmod d\). Choose the common residue \(r\) with \(0\le r<d\). Since \(a_2,b_3\ge0\), the exponents \((a_2-r)/d\) and \((b_3-r)/d\) are nonnegative, and prime-by-prime division produces integers \(C,D\) satisfying \eqref{eq:rootfactor}. The converse follows immediately by taking valuations. The Kummer formulation is the same statement modulo \(d\)-th powers.

Finally,
\[
  2^x=A=2^rC^d=2^{r+d\log C/\log2},
\]
so injectivity of the real exponential gives \(x=r+d\log C/\log2\). The analogous identity for base \(3\) yields the same \(y\), proving \eqref{eq:descentexp}.
\end{proof}

\begin{corollary}[One-step saturation descent]\label{cor:onestep}
Taking \(d=\defect(A,B)\) in \cref{thm:rootextract} produces a counterexample pair \((C,D)\) with
\[
       \defect(C,D)=1.
\]
\end{corollary}

\begin{proof}
After subtracting the common residue and dividing all relevant valuations by the full gcd \(\defect(A,B)\), the gcd of the new collection is \(1\).
\end{proof}

Thus there is no need for an indefinite descent merely to reach a saturated pair: the full defect removes every common root at once.

\section{Canonical primitive normalization}

The saturation defect handles common roots, while a separate integer shift handles common powers of the distinguished axis bases.

\begin{definition}[Axis content]\label{def:axiscontent}
For a counterexample pair \((A,B)\), define
\[
       c(A,B)=\min\bigl(v_2(A),v_3(B)\bigr).
\]
Then
\[
 A'=A/2^{c(A,B)},\qquad B'=B/3^{c(A,B)}
\]
are positive integers and correspond to the shifted exponent \(x'=x-c(A,B)\).
\end{definition}

The difference \(v_3(B)-v_2(A)\) and all off-axis valuations are unchanged by a matched axis shift, so
\[
       \defect(A',B')=\defect(A,B).
\]

\begin{definition}[Canonical primitive counterexample]\label{def:canonical}
A nonintegral common exponent \(\beta\), with
\[
        A_0=2^\beta,\qquad B_0=3^\beta,
\]
is called \emph{canonical primitive} if
\begin{equation}\label{eq:canonicalconditions}
  \defect(A_0,B_0)=1,
  \qquad
  \min\bigl(v_2(A_0),v_3(B_0)\bigr)=0.
\end{equation}
\end{definition}

\begin{proposition}[Existence of a canonical primitive root]\label{prop:canonicalexist}
If the Alaoglu--Erd\H{o}s conjecture is false, then a canonical primitive counterexample exists.
\end{proposition}

\begin{proof}
Start from any counterexample \(x\). Apply \cref{cor:onestep} to extract the full defect, obtaining a saturated counterexample \(y\). Then subtract its axis content. The resulting exponent is still nonintegral, because only an integer was subtracted; its two powers remain positive integers; its defect remains \(1\); and its axis content is \(0\). It cannot be negative, and it cannot be \(0\), so \cref{lem:elementary} gives \(\beta\ge1\).
\end{proof}

From now on, whenever a canonical primitive counterexample is under discussion, write
\begin{equation}\label{eq:canonicalpair}
       A=2^\beta,\qquad B=3^\beta,
\end{equation}
with
\begin{equation}\label{eq:canonicalshort}
       \defect(A,B)=1,
       \qquad
       \min(a_2,b_3)=0.
\end{equation}

\section{The conditional free-monoid theorem}

The next theorem is the strongest structural result of this report. It turns the vague phrase ``all further solutions lie in the same rank-two plane'' into exact group and monoid equalities.

\begin{theorem}[Common rational exponent group]\label{thm:rationalgroup}
Assume \(\beta\) is canonical primitive. Then
\begin{equation}\label{eq:Erat}
       \Erat(2,3)=\Z+\Z\beta.
\end{equation}
Every element has a unique representation \(n+k\beta\) with \(n,k\in\Z\).
\end{theorem}

\begin{proof}
By \cref{cor:Ealgplane}, every \(t\in\Erat(2,3)\) has the form \(t=q+r\beta\) with \(q,r\in\Q\). Its valuation vector is
\[
     \val(2^t,3^t)=qg_0+rg_1.
\]
Because both values are rational, this vector is integral, hence belongs to \(L(A,B)^{\sat}\). The canonical condition \(\defect(A,B)=1\) and \cref{thm:satindex} say that \(L(A,B)\) is saturated. Therefore \(qg_0+rg_1\in L(A,B)\). Linear independence of \(g_0,g_1\) gives \(q,r\in\Z\). The reverse inclusion follows from
\[
  \Pair(n+k\beta)=\axis^n(A,B)^k\in\Q_{>0}^2.
\]
Uniqueness follows from the irrationality of \(\beta\).
\end{proof}

\begin{theorem}[Conditional free-monoid structure]\label{thm:freemonoid}
Assume \(\beta\) is canonical primitive. Then
\begin{equation}\label{eq:Eintmonoid}
       \Eint(2,3)=\N_0+\N_0\beta.
\end{equation}
In particular, the common integral exponents form a free commutative monoid on the two generators \(1\) and \(\beta\).
\end{theorem}

\begin{proof}
Let \(t\in\Eint(2,3)\). By \cref{thm:rationalgroup}, write \(t=n+k\beta\) with \(n,k\in\Z\). Since \(\beta\notin\Z\), the integer \(A=2^\beta\) is not a pure power of \(2\). Hence some prime \(p\ne2\) has \(a_p>0\). At that prime,
\[
       v_p(2^t)=k a_p\ge0,
\]
so \(k\ge0\). At the distinguished primes,
\[
       v_2(2^t)=n+k a_2\ge0,
       \qquad
       v_3(3^t)=n+k b_3\ge0.
\]
Because \(\min(a_2,b_3)=0\), one of these inequalities is simply \(n\ge0\). Thus \(n,k\in\N_0\).

Conversely, if \(n,k\in\N_0\), then
\[
       2^{n+k\beta}=2^nA^k,
       \qquad
       3^{n+k\beta}=3^nB^k
\]
are positive integers.
\end{proof}

\begin{corollary}[Uniqueness and least counterexample]\label{cor:uniqueprimitive}
If the conjecture is false, there is exactly one canonical primitive counterexample \(\beta\). It is the least nonintegral element of \(\Eint(2,3)\), and every nonintegral element is uniquely
\[
       n+k\beta,\qquad n\in\N_0,\quad k\in\N.
\]
\end{corollary}

\begin{proof}
For a canonical \(\beta\), \cref{thm:freemonoid} shows that every nonintegral solution is at least \(\beta\). If \(\beta'\) were another canonical primitive counterexample, then \(\beta'\ge\beta\) by the description attached to \(\beta\), and symmetrically \(\beta\ge\beta'\). Hence they are equal. The representation is unique because \(\beta\) is irrational.
\end{proof}

\begin{remark}[What the theorem does and does not prove]
The free-monoid theorem is conditional: it describes the entire solution set if one counterexample exists. It does not exclude the free monoid. Its value is that every future proof attempt may work against the sharply normalized conditions \eqref{eq:canonicalshort}, and any operation producing a common exponent outside \(\N_0+\N_0\beta\) gives an immediate contradiction.
\end{remark}

\section{All descendants: shift, power, and exact defect}

The free-monoid theorem gives a canonical coordinate system. Let
\[
       \alpha=n+k\beta,
       \qquad n\in\N_0,\quad k\in\N,
\]
and write
\[
       C=2^\alpha=2^nA^k,
       \qquad
       D=3^\alpha=3^nB^k.
\]

\begin{theorem}[Orbit descendant formulas]\label{thm:descendant}
For every nonintegral descendant \(\alpha=n+k\beta\),
\begin{align}
  c(C,D)&=n,\label{eq:descaxis}\\
  \defect(C,D)&=k,\label{eq:descdefect}\\
  \frac{(C,D)}{(2,3)^n}&=(A^k,B^k).\label{eq:localizedpower}
\end{align}
Consequently the canonical reduction procedure applied to \((C,D)\) returns \((A,B)\) and \(\beta\).
\end{theorem}

\begin{proof}
At the distinguished primes,
\[
 v_2(C)=n+k a_2,
 \qquad
 v_3(D)=n+k b_3.
\]
Since \(\min(a_2,b_3)=0\), their minimum is \(n\), proving \eqref{eq:descaxis}. For primes away from the axes, the valuations are \(k a_p\) and \(k b_p\), while
\[
 v_3(D)-v_2(C)=k(b_3-a_2).
\]
Their gcd is \(k\defect(A,B)=k\), proving \eqref{eq:descdefect}. Equation \eqref{eq:localizedpower} is immediate.

After stripping the axis factor, the pair is the coordinatewise \(k\)-th power of \((A,B)\); extracting its full defect \(k\) recovers \((A,B)\). Any residue introduced by reducing \(n\) modulo \(k\) is removed by the subsequent axis normalization, so the canonical output is independent of the order in which the two elementary normalizations are described.
\end{proof}

\begin{corollary}[Intrinsic meaning of the second coordinate]\label{cor:kintrinsic}
For a nonintegral common integral exponent \(\alpha\), the coefficient \(k\) in \(\alpha=n+k\beta\) is intrinsic: it is the saturation index of the valuation lattice generated by \((2,3)\) and \((2^\alpha,3^\alpha)\).
\end{corollary}

This provides a useful diagnostic. A candidate pair with defect \(k>1\) can never be the primitive counterexample; it is necessarily a \(k\)-fold inflation of the unique primitive orbit, after an axis shift.

\subsection{Integer-shift invariance}

Replacing \(\beta\) by \(\beta+h\), \(h\in\Z\), changes the pair to
\[
       (2^hA,3^hB),
\]
but leaves the rational exponent group \(\Z+\Z\beta\) unchanged. On coefficient vectors, the change is the unimodular transformation
\[
       n+k\beta=(n-hk)+k(\beta+h).
\]
The off-axis valuations and the difference \(b_3-a_2\) are unchanged, so the saturation defect is shift-invariant. The axis normalization in \cref{def:canonical} selects one preferred representative of this integer-shift orbit.

\section{Smith normal form of the determinant lift}

The index \(k\) in \cref{thm:descendant} also appears twice in a natural rank-three lattice. This is the exact Smith-normal-form statement announced in the preliminary analysis.

Let \(M\cong\Z^2\) be an exponent-coordinate lattice. An integer matrix \(G\in M_2(\Z)\) acts on \(M\), and it acts on the rank-one orientation lattice \(\bigwedge^2M\) by multiplication by \(\det G\). Define the \emph{determinant lift}
\[
       \widehat G=G\oplus(\det G)
       :M\oplus\bigwedge^2M\longrightarrow
        M\oplus\bigwedge^2M.
\]
In the standard basis,
\[
  G=\begin{pmatrix}a&b\\c&d\end{pmatrix},
  \qquad
  \widehat G=
  \begin{pmatrix}
     a&b&0\\
     c&d&0\\
     0&0&\Delta
  \end{pmatrix},
  \qquad \Delta=ad-bc.
\]

\begin{theorem}[Exact Smith form of the determinant lift]\label{thm:snflift}
Suppose \(\gcd(a,b,c,d)=1\) and \(\Delta\ne0\). Then
\begin{equation}\label{eq:snflift}
       \SNF(\widehat G)=
       \operatorname{diag}(1,|\Delta|,|\Delta|).
\end{equation}
Equivalently,
\[
  \operatorname{coker}(\widehat G)
       \cong (\Z/|\Delta|\Z)^2.
\]
\end{theorem}

\begin{proof}
The gcd of all entries of \(\widehat G\) is \(1\), so the first determinantal divisor is \(1\). Every nonzero \(2\times2\) minor is either \(\Delta\) or \(\Delta\) times one of \(a,b,c,d\). Since the upper-left minor is exactly \(\Delta\), the gcd of all \(2\times2\) minors is \(|\Delta|\). Finally,
\[
       |\det\widehat G|=|\Delta|^2.
\]
The invariant factors are therefore \(1,|\Delta|,|\Delta|\).
\end{proof}

For a descendant \(\alpha=n+k\beta\), the basis change from \((1,\beta)\) to \((1,\alpha)\) is
\[
       G_{n,k}=
       \begin{pmatrix}1&n\\0&k\end{pmatrix},
       \qquad \det G_{n,k}=k.
\]

\begin{corollary}[Descendant Smith form]\label{cor:descSNF}
For \(k\ge1\),
\[
 \SNF\!\begin{pmatrix}
       1&n&0\\
       0&k&0\\
       0&0&k
 \end{pmatrix}
 =\operatorname{diag}(1,k,k).
\]
\end{corollary}

The two copies of \(k\) should not be mistaken for an additional transcendence theorem. One copy records the index of the exponent sublattice; the second records the induced index in the orientation coordinate. What is useful is that there is no hidden third invariant: the rank-three bookkeeping is completely controlled by the same scalar \(k\) that appears in the perfect-power descent.

\section{One additional iterate would finish the conjecture}

The original hypothesis gives two algebraic rows and two algebraic columns:
\[
 \begin{pmatrix}
  2&2^x\\
  3&3^x
 \end{pmatrix}.
\]
Six Exponentials needs a third exponent. The cleanest possible third column is \(x^2\). The next theorem shows that this single orbit extension is already decisive.

\begin{theorem}[Three-step algebraic orbit theorem]\label{thm:threeorbit}
Let \(a,b>0\) be multiplicatively independent algebraic numbers, neither equal to \(1\). Let \(\xi\in\R\). If all six numbers
\[
       a,\ a^\xi,\ a^{\xi^2},
       \qquad
       b,\ b^\xi,\ b^{\xi^2}
\]
are algebraic, then \(\xi\in\Q\).
\end{theorem}

\begin{proof}
Assume \(\xi\notin\Q\). Since \(a\) and \(a^\xi\) are algebraic and \(a\notin\{0,1\}\), Gelfond--Schneider shows that \(\xi\) cannot be algebraic irrational. Hence \(\xi\) is transcendental, and therefore \(1,\xi,\xi^2\) are linearly independent over \(\Q\). Multiplicative independence of \(a,b\) makes \(\log a,\log b\) linearly independent over \(\Q\). Applying Six Exponentials to these two logarithms and the three exponents contradicts the algebraicity of all six displayed numbers.
\end{proof}

\begin{corollary}[One-step orbit barrier]\label{cor:onestepbarrier}
If \(\beta\) is a hypothetical nonintegral solution and both
\[
       2^{\beta^2}=A^\beta,
       \qquad
       3^{\beta^2}=B^\beta
\]
are algebraic, then \(\beta\in\Q\), hence \(\beta\in\Z\). Therefore a proof of algebraicity of just the next simultaneous iterate would prove the Alaoglu--Erd\H{o}s conjecture.
\end{corollary}

\begin{warning}
The original assumptions give \(A,B\in\N\), but they do not imply \(A^\beta,B^\beta\in\Qbar\). Treating the latter as automatic would simply insert the missing theorem into the argument. The value of \cref{cor:onestepbarrier} is to identify the exact missing closure statement.
\end{warning}

\section{Rational Hankel triples and the square valuation pattern}

The three-step orbit can be encoded by a rank-one logarithmic Hankel matrix. For positive \(U,V,W\) with \(U\ne1\), define
\[
      H(U,V,W)=
      \begin{pmatrix}
       \log U&\log V\\
       \log V&\log W
      \end{pmatrix}.
\]

\begin{lemma}[Scalar Hankel criterion]\label{lem:hankelscalar}
For positive \(U,V,W\) with \(U\ne1\), the identity
\begin{equation}\label{eq:hankelscalar}
       (\log V)^2=\log U\,\log W
\end{equation}
holds if and only if, with
\[
       \xi=\frac{\log V}{\log U},
\]
one has
\[
       V=U^\xi,
       \qquad
       W=U^{\xi^2}.
\]
\end{lemma}

\begin{proof}
The first equality is the definition of \(\xi\). Equation \eqref{eq:hankelscalar} gives
\[
       \log W=\frac{(\log V)^2}{\log U}=\xi^2\log U,
\]
and exponentiation gives the second. The converse is immediate.
\end{proof}

A single rational Hankel triple can have an unknown irrational slope; the second multiplicatively independent base is what activates Six Exponentials.

\begin{theorem}[Two-base rational Hankel theorem]\label{thm:twohankel}
For \(i=1,2\), let \(U_i,V_i,W_i\) be positive algebraic numbers with \(U_i\ne1\). Suppose
\begin{equation}\label{eq:twoslope}
  \frac{\log V_1}{\log U_1}
   =\frac{\log V_2}{\log U_2}=\xi
\end{equation}
and each triple satisfies \eqref{eq:hankelscalar}. If \(U_1,U_2\) are multiplicatively independent, then \(\xi\in\Q\).
\end{theorem}

\begin{proof}
By \cref{lem:hankelscalar}, the six entries are
\[
 U_1,U_1^\xi,U_1^{\xi^2},
 \qquad
 U_2,U_2^\xi,U_2^{\xi^2}.
\]
Apply \cref{thm:threeorbit}.
\end{proof}

The rational slope has a stronger integrality consequence than the first-step perfect-power condition alone.

\begin{theorem}[Perfect-square valuation normal form]\label{thm:squarepattern}
Let \(U,V,W\in\Q_{>0}\), let \(\xi=r/s\in\Q\) be in lowest terms with \(s>0\), and suppose
\[
       V=U^\xi,
       \qquad
       W=U^{\xi^2}.
\]
Then there is a unique \(T\in\Q_{>0}\) such that
\begin{equation}\label{eq:squarepower}
       U=T^{s^2},
       \qquad
       V=T^{rs},
       \qquad
       W=T^{r^2}.
\end{equation}
Equivalently, for every prime \(p\),
\begin{equation}\label{eq:squarevaluations}
  \bigl(v_p(U),v_p(V),v_p(W)\bigr)
      =d_p(s^2,rs,r^2)
\end{equation}
for an integer \(d_p\), with only finitely many nonzero \(d_p\). The triple obeys the two monomial recurrences
\begin{equation}\label{eq:monomialrec}
       U^r=V^s,
       \qquad
       V^r=W^s.
\end{equation}
Conversely, \eqref{eq:squarepower} implies the orbit identities, the Hankel identity, and \eqref{eq:monomialrec}.
\end{theorem}

\begin{proof}
Let \(u_p=v_p(U)\), \(v_p=v_p(V)\), and \(w_p=v_p(W)\). From \(V=U^{r/s}\),
\[
       s v_p=r u_p.
\]
Coprimality gives \(s\mid u_p\), so write \(u_p=s e_p\) and \(v_p=r e_p\). From \(W=V^{r/s}\),
\[
       s w_p=r v_p=r^2e_p.
\]
Because \(\gcd(r,s)=1\), we have \(s\mid e_p\). Write \(e_p=s d_p\). Then
\[
       (u_p,v_p,w_p)=d_p(s^2,rs,r^2).
\]
Taking \(T=\prod_p p^{d_p}\) proves \eqref{eq:squarepower}; uniqueness follows from valuations. The recurrences and converse are immediate by comparing exponents.
\end{proof}

\begin{corollary}[Integral version]\label{cor:integralsquare}
If \(U,V,W\in\N\) and \(r,s>0\), then the \(T\) in \cref{thm:squarepattern} lies in \(\N\). Thus both endpoints are perfect squares whenever \(r,s>1\), and more precisely they are an \(s^2\)-th and an \(r^2\)-th power of the same primitive base.
\end{corollary}

The exponent vector \((s^2,rs,r^2)\) is also the primitive kernel direction of the two recurrence rows
\[
  R_{r,s}=
  \begin{pmatrix}
    r&-s&0\\
    0&r&-s
  \end{pmatrix}.
\]
When \(\gcd(r,s)=1\), the gcd of the \(2\times2\) minors \(r^2,-rs,s^2\) is \(1\), so the row lattice is saturated and its integer kernel is exactly
\[
        \Z(s^2,rs,r^2).
\]
This is the elementary lattice source of the square pattern.

\section{The exact rank-three branch inside the canonical orbit}

The preceding theorem treats an abstract exponentiation orbit. Under a canonical counterexample, every rational point is already coordinatized by the lattice \(\Z+\Z\beta\), allowing a second, purely algebraic classification.

Let
\[
       t_i=n_i+k_i\beta\in\Erat(2,3),
       \qquad i=0,1,2,
\]
and define
\[
       U_i=2^{t_i},\qquad V_i=3^{t_i}.
\]
Because \(\beta\) is transcendental, equality at \(\beta\) between polynomials with rational coefficients is a polynomial identity.

\begin{theorem}[Global Hankel-ray classification]\label{thm:globalhankel}
Assume \(t_0,t_1,t_2\ne0\) and
\begin{equation}\label{eq:texphankel}
       t_1^2=t_0t_2.
\end{equation}
Then there is a primitive vector \((a,b)\in\Z^2\), integers \(m_0,m_1,m_2\), and
\[
       \tau=a+b\beta
\]
such that
\[
       t_i=m_i\tau,
       \qquad
       m_1^2=m_0m_2.
\]
If the \(t_i\) belong to \(\Eint(2,3)\) and are positive, then one may choose \(a,b,m_i\in\N_0\), with \(\gcd(a,b)=1\). If in addition all \(m_i>0\), there exist \(d,r,s\in\N\), \(\gcd(r,s)=1\), such that
\begin{equation}\label{eq:globalray}
       (t_0,t_1,t_2)
        =d(s^2,rs,r^2)\tau.
\end{equation}
\end{theorem}

\begin{proof}
Set \(L_i(X)=n_i+k_iX\in\Z[X]\). Equation \eqref{eq:texphankel} gives
\[
       L_1(\beta)^2=L_0(\beta)L_2(\beta).
\]
Transcendence of \(\beta\) implies
\[
       L_1(X)^2=L_0(X)L_2(X)
       \quad\text{in }\Q[X].
\]
If all \(L_i\) are constant, the conclusion is immediate. Otherwise, the product on the right has degree two, so all three polynomials are nonconstant linear polynomials. Unique factorization in \(\Q[X]\) forces them to be scalar multiples of one linear polynomial. Choose its integer coefficient vector \((a,b)\) primitive. Gauss's lemma then makes the scalar multipliers integers, giving \(t_i=m_i\tau\) and \(m_1^2=m_0m_2\).

For positive integral exponents, \cref{thm:freemonoid} puts all coefficient vectors in \(\N_0^2\). A common primitive ray can therefore be oriented with \(a,b\ge0\), and the multipliers are nonnegative. Finally, the standard parametrization of positive integer solutions to \(m_1^2=m_0m_2\) gives
\[
       (m_0,m_1,m_2)=d(s^2,rs,r^2)
\]
with \(\gcd(r,s)=1\).
\end{proof}

\begin{corollary}[Pair-valued rank-three normal form]\label{cor:pairhankel}
Under the positive integral hypotheses of \cref{thm:globalhankel}, put
\[
       \mathbf T=\Pair(\tau)^d.
\]
Then, componentwise,
\begin{equation}\label{eq:pairpattern}
       \bigl(\Pair(t_0),\Pair(t_1),\Pair(t_2)\bigr)
       =\bigl(\mathbf T^{s^2},\mathbf T^{rs},\mathbf T^{r^2}\bigr),
\end{equation}
and
\[
       \Pair(t_0)^r=\Pair(t_1)^s,
       \qquad
       \Pair(t_1)^r=\Pair(t_2)^s.
\]
For every prime and each of the two coordinates, the three valuations are an integer multiple of \((s^2,rs,r^2)\).
\end{corollary}

\subsection{Localized version}

Write each solution pair uniquely as
\[
       \Pair(n+k\beta)=\axis^n(A,B)^k.
\]
After quotienting by the axis subgroup generated by \(\axis\), only the coefficient \(k\) remains. A non-axis localized Hankel triple is therefore determined by integers \(k_0,k_1,k_2>0\) with
\[
       k_1^2=k_0k_2.
\]
It follows that
\[
       (k_0,k_1,k_2)=d(s^2,rs,r^2),
\]
and the localized pair triple is
\[
  \bigl((A,B)^{k_0},(A,B)^{k_1},(A,B)^{k_2}\bigr)
   =\bigl(\mathbf T^{s^2},\mathbf T^{rs},\mathbf T^{r^2}\bigr),
  \qquad \mathbf T=(A,B)^d.
\]
This is the exact ``one primitive orbit'' statement: after the known axis is removed, every rank-three Hankel branch is a square-pattern power of one localized generator.

\begin{remark}[Exact converse]
Given a positive rational pair \(\mathbf T\), coprime \(r,s\in\N\), and \(d\in\N\), the three pairs in \eqref{eq:pairpattern} satisfy both monomial recurrences and the scalar logarithmic Hankel equation in each coordinate. Conversely, a positive localized rational triple satisfying the two recurrences has precisely this form. The scalar logarithmic equation alone identifies a real orbit slope; the monomial recurrences are what certify that the slope is rational and that the valuation vector is integral.
\end{remark}

\section{Fractional-linear closure and why it does not escape}

The rank-two lattice naturally generates many new-looking exponents when the base pair is changed. These are useful for organizing the problem, but they do not create a third independent direction.

For nonzero \(t,u\in\Erat(2,3)\), set
\[
       \theta=\frac{u}{t}.
\]
Then
\[
       \Pair(t)^\theta=\Pair(u).
\]
Under the canonical coordinates,
\[
 t=n+k\beta,
 \qquad
 u=m+\ell\beta,
\]
so
\begin{equation}\label{eq:mobius}
       \theta=\frac{m+\ell\beta}{n+k\beta}.
\end{equation}

\begin{proposition}[Transcendence of nondegenerate orbit ratios]\label{prop:mobtrans}
If the coefficient vectors \((n,k)\) and \((m,\ell)\) are not proportional over \(\Q\), then the number \(\theta\) in \eqref{eq:mobius} is transcendental.
\end{proposition}

\begin{proof}
The ratio cannot be rational: cross multiplication would give a nontrivial linear relation between \(1\) and \(\beta\). If it were algebraic irrational, then a positive rational coordinate of \(\Pair(t)\), raised to the algebraic irrational power \(\theta\), would be the corresponding rational coordinate of \(\Pair(u)\), contradicting Gelfond--Schneider. Therefore it is transcendental.
\end{proof}

\begin{proposition}[Common exponent group after changing bases]\label{prop:changedbasegroup}
For nonzero \(t\in\Erat(2,3)\),
\begin{equation}\label{eq:changedbase}
  \Erat\bigl(2^t,3^t\bigr)
   =\frac{1}{t}\bigl(\Z+\Z\beta\bigr).
\end{equation}
\end{proposition}

\begin{proof}
The pair \(((2^t)^\theta,(3^t)^\theta)\) is rational exactly when \(\Pair(t\theta)\) is rational. Apply \cref{thm:rationalgroup}.
\end{proof}

Thus every base pair arising from the orbit again has a rank-two common exponent group. The following statement gives an elementary internal version of the three-step obstruction.

\begin{proposition}[Quadratic closure inside the orbit is rational]\label{prop:quadraticclosure}
Let \(t\ne0\) belong to \(\Erat(2,3)\), and let \(\theta\in\Erat(2^t,3^t)\). If also \(\theta^2\in\Erat(2^t,3^t)\), then \(\theta\in\Q\).
\end{proposition}

\begin{proof}
There exist \(u,v\in\Z+\Z\beta\) with \(u=t\theta\) and \(v=t\theta^2\). Hence
\[
       u^2=tv.
\]
Write \(t,u,v\) as linear polynomials in \(\beta\). Transcendence of \(\beta\) turns this equality into a polynomial identity. The unique-factorization argument in \cref{thm:globalhankel} forces the coefficient vectors of \(t,u,v\) onto one rational ray. Therefore \(u/t=\theta\) is rational.
\end{proof}

This proposition explains why fractional-linear manipulations alone cannot manufacture the missing \(\beta^2\) column. They move among ratios of linear forms in \(\beta\), and quadratic closure collapses to a rational slope.

\section{Reciprocal duality for the primitive pair}

The canonical pair \((A,B)=(2^\beta,3^\beta)\) reverses the roles of \(\beta\) and \(1/\beta\). Indeed,
\[
       A^{1/\beta}=2,
       \qquad
       B^{1/\beta}=3.
\]
The free-monoid theorem makes this reciprocity exact at the level of all common exponents.

\begin{theorem}[Reciprocal group and monoid]\label{thm:reciprocal}
For the canonical primitive pair \((A,B)\),
\begin{align}
 \Erat(A,B)&=\Z+\frac{1}{\beta}\Z,\label{eq:reciprat}\\
 \Eint(A,B)&=\N_0+\frac{1}{\beta}\N_0.\label{eq:recipint}
\end{align}
\end{theorem}

\begin{proof}
For any real \(t\),
\[
       (A^t,B^t)=\Pair(\beta t).
\]
The pair is rational exactly when \(\beta t\in\Z+\Z\beta\), which is equivalent to \(t\in\Z+\Z/\beta\). It is integral exactly when \(\beta t\in\N_0+\N_0\beta\), giving \eqref{eq:recipint}.
\end{proof}

\begin{remark}
The reciprocal exponent \(1/\beta\) is a genuine common integral exponent for the \emph{new} bases \((A,B)\), not for the original bases \((2,3)\). This distinction prevents an invalid attempt to adjoin \(1/\beta\) as a third exponent in \(\Ealg(2,3)\). The duality is exact but remains two-dimensional.
\end{remark}

More generally, for any positive \(t=n+k\beta\), the common rational exponent group of \((2^t,3^t)\) is the fractional-linear lattice \((\Z+\Z\beta)/t\). This family is closed under composition: an exponent taking \(\Pair(t)\) to \(\Pair(u)\), followed by one taking \(\Pair(u)\) to \(\Pair(v)\), multiplies to \(v/t\). The resulting groupoid is a useful exact model of all transformations obtainable without adding new arithmetic information.

\section{Continued-fraction localization at the axis}

The canonical normalization has a useful consequence for Diophantine approximations to \(\beta\). Let \(p/q\) be a positive rational approximation, with \(q\ge1\), and put
\[
       \varepsilon=q\beta-p.
\]
Then
\begin{equation}\label{eq:nearaxis}
 \Pair(\varepsilon)=
 \left(\frac{A^q}{2^p},\frac{B^q}{3^p}\right).
\end{equation}
For continued-fraction convergents \(p_j/q_j\), the exponents \(\varepsilon_j\) tend to zero and the rational pairs in \eqref{eq:nearaxis} tend to \((1,1)\).

\begin{proposition}[Exact denominator dichotomy]\label{prop:denomdichotomy}
For a canonical primitive pair, at least one of the following holds for every \(p,q\ge1\):
\begin{enumerate}[label=\textup{(\roman*)}]
\item \(a_2=0\), in which case \(A^q/2^p\) is in lowest terms and has denominator exactly \(2^p\);
\item \(b_3=0\), in which case \(B^q/3^p\) is in lowest terms and has denominator exactly \(3^p\).
\end{enumerate}
Both may hold.
\end{proposition}

\begin{proof}
The canonical condition gives \(\min(a_2,b_3)=0\). If \(a_2=0\), then \(A\) is odd, so \(\gcd(A^q,2^p)=1\). The base-three statement is identical.
\end{proof}

This prevents simultaneous cancellation at the axis. It also gives a completely elementary, though weak, lower bound for the approximation error.

\begin{corollary}[Elementary near-axis lower bound]\label{cor:nearaxisbound}
Assume \(|\varepsilon|\le1\). If \(a_2=0\), then
\begin{equation}\label{eq:eps2bound}
       |\varepsilon|\ge
       \frac{1}{2^{p+1}\log2}.
\end{equation}
If \(b_3=0\), then
\begin{equation}\label{eq:eps3bound}
       |\varepsilon|\ge
       \frac{1}{3^{p+1}\log3}.
\end{equation}
\end{corollary}

\begin{proof}
Suppose \(a_2=0\). The reduced rational \(A^q/2^p\ne1\), so
\[
 \left|\frac{A^q}{2^p}-1\right|\ge2^{-p}.
\]
On the other hand, by the mean-value theorem,
\[
       |2^\varepsilon-1|
       \le 2\log2\,|\varepsilon|
       \qquad (|\varepsilon|\le1).
\]
Combine this with \eqref{eq:nearaxis}. The proof for base \(3\) is the same.
\end{proof}

For convergents, \(|\varepsilon_j|<1/q_{j+1}\), so the preceding inequalities bound \(q_{j+1}\) by an exponential function of \(p_j\). Baker's theory of linear forms in logarithms gives much stronger effective lower bounds for
\[
       |q\log A-p\log2|=|\varepsilon|\log2
\]
when \(A\) is fixed. Nevertheless, both bounds address a \emph{single} linear form. The exact relation defining a counterexample is the vanishing of a determinant of logarithms,
\[
       \log A\log3-\log B\log2=0,
\]
and the two linear forms attached to \((A,2)\) and \((B,3)\) are exactly proportional. Merely placing separate Baker bounds on them does not create a contradiction.

\subsection{Synchronized difference pairs}

For \(\varepsilon=q\beta-p<0\), define positive integers
\[
       \Delta_2=2^p-A^q,
       \qquad
       \Delta_3=3^p-B^q.
\]
Then
\[
 \frac{\Delta_2}{2^p}=1-2^\varepsilon,
 \qquad
 \frac{\Delta_3}{3^p}=1-3^\varepsilon.
\]
As \(\varepsilon\to0\),
\begin{equation}\label{eq:syncasymp}
 \frac{\Delta_3/3^p}{\Delta_2/2^p}
   =\frac{\log3}{\log2}\bigl(1+O(|\varepsilon|)\bigr).
\end{equation}
The formula is rigorous after replacing the \(O\)-term by the standard exponential Taylor remainder. It shows that every counterexample produces synchronized near-collisions of two independent exponential sequences. A promising arithmetic route would need information coupling the prime divisors of \(\Delta_2\) and \(\Delta_3\), not merely separate size estimates. No such coupling is presently proved here.

\section{Local Kummer form of the defect}

The saturation defect has a useful prime-by-prime reformulation. Let \(\ell\) be a prime, and reduce the two valuation columns \(g_0,g_1\) modulo powers of \(\ell\).

\begin{theorem}[Local defect exponent]\label{thm:localdefect}
For every prime \(\ell\),
\begin{equation}\label{eq:localdefect}
 v_\ell\bigl(\defect(A,B)\bigr)
 =\max\left\{e\ge0:
    g_1\equiv r g_0\pmod{\ell^e}
    \text{ for some }r\in\Z/\ell^e\Z
   \right\}.
\end{equation}
Equivalently, \(\ell^e\mid\defect(A,B)\) exactly when there is a common residue \(r\pmod{\ell^e}\) such that
\begin{equation}\label{eq:localfactor}
       A=2^r C^{\ell^e},
       \qquad
       B=3^r D^{\ell^e}
\end{equation}
for positive integers \(C,D\).
\end{theorem}

\begin{proof}
Because \(g_0\) has a unit entry in each distinguished coordinate, congruence of \(g_1\) to a scalar multiple of \(g_0\) says precisely that all off-axis valuations are divisible by \(\ell^e\) and that \(a_2\equiv b_3\pmod{\ell^e}\). This is equivalent to divisibility of every generator in \eqref{eq:defect}, and \cref{thm:rootextract} gives the factorization.
\end{proof}

\begin{corollary}[Canonical mod-\(\ell\) independence]\label{cor:modlind}
For a canonical primitive pair and every prime \(\ell\), the two valuation columns \(g_0,g_1\) are linearly independent over \(\mathbb F_\ell\). Equivalently, for every \(\ell\), at least one of the following holds:
\[
 \ell\nmid a_p\text{ for some }p\ne2,
 \quad
 \ell\nmid b_p\text{ for some }p\ne3,
 \quad
 a_2\not\equiv b_3\pmod\ell.
\]
\end{corollary}

This is an exact local-global target. To prove the conjecture by descent, it would suffice to show that a hypothetical archimedean relation
\[
       \frac{\log A}{\log2}=\frac{\log B}{\log3}\notin\Z
\]
forces the two valuation columns to become dependent modulo at least one prime \(\ell\). The canonical theorem says that a counterexample must evade this dependence for \emph{every} \(\ell\).

\section{Sharp sufficient criteria for a complete proof}

The structural results isolate several independent ways to finish the conjecture. Each criterion below is unconditional as an implication; what remains open is its hypothesis.

\begin{criterion}[Defect forcing]\label{crit:defect}
Suppose one proves that every nonintegral solution pair \((A,B)\) satisfies \(\defect(A,B)>1\). Then the Alaoglu--Erd\H{o}s conjecture follows.
\end{criterion}

\begin{proof}
A false conjecture would have a canonical primitive pair with defect \(1\), contradicting the assumed forcing statement.
\end{proof}

A prime-local variant is enough: it suffices to prove that every counterexample has some prime \(\ell\) for which \(g_1\) is an axis multiple of \(g_0\) modulo \(\ell\).

\begin{criterion}[Nonlinear common exponent]\label{crit:nonlinear}
Suppose \(\beta\) is a hypothetical counterexample and one can construct
\[
       \gamma\in\Ealg(2,3)\setminus(\Q+\Q\beta).
\]
Then the conjecture follows.
\end{criterion}

\begin{proof}
This contradicts \cref{cor:Ealgplane}. Equivalently, \(1,\beta,\gamma\) would supply the three independent exponent columns required by Six Exponentials.
\end{proof}

\begin{corollary}[Polynomial closure criterion]\label{cor:polyclosure}
It is enough to prove that, for some nonaffine polynomial \(F\in\Q[X]\), both \(2^{F(\beta)}\) and \(3^{F(\beta)}\) are algebraic. In particular, \(F(X)=X^2\) gives \cref{cor:onestepbarrier}.
\end{corollary}

\begin{proof}
Since \(\beta\) is transcendental, a nonaffine polynomial \(F(\beta)\) cannot lie in \(\Q+\Q\beta\) unless \(F\) becomes affine after cancellation, which a polynomial does not.
\end{proof}

\begin{criterion}[Monoid leakage]\label{crit:leakage}
For a canonical primitive \(\beta\), any proof that \(\Eint(2,3)\) contains an element \(n+k\beta\) with \(n<0\) or \(k<0\) yields a contradiction.
\end{criterion}

This criterion packages all possible subtraction or reciprocal-closure strategies. The original solution set is closed under addition, but not known to be closed under subtraction. Any genuinely arithmetic mechanism that creates one forbidden negative coefficient would finish the problem.

\begin{criterion}[Matched root extraction]\label{crit:root}
If a counterexample pair necessarily admits, for some \(d\ge2\), a factorization
\[
       A=2^rC^d,
       \qquad
       B=3^rD^d
\]
with the same \(r\), then the conjecture follows.
\end{criterion}

This is the concrete perfect-power form of \cref{crit:defect}. It is often more approachable than the gcd definition because it can be tested in Kummer quotients.

\begin{criterion}[One more simultaneous iterate]\label{crit:iterate}
If the original integrality of \(2^\beta,3^\beta\) can be shown to imply algebraicity of \(2^{\beta^2},3^{\beta^2}\), then the conjecture follows.
\end{criterion}

\begin{criterion}[A third orbit ray]\label{crit:thirdray}
If there exists any positive algebraic base pair \((u,v)\) with \(\log u,\log v\) independent over \(\Q\) and three independent common algebraic exponents constructed from a hypothetical \(\beta\), then Six Exponentials gives a contradiction.
\end{criterion}

The fractional-linear orbit analysis in \cref{prop:changedbasegroup} shows that changing from \((2,3)\) to a pair already in the same orbit does not meet this criterion: its common exponent group remains two-dimensional.

\section{What the new reductions rule out}

A productive progress report should identify not only successful lemmas but also tempting arguments that the lemmas show to be circular or insufficient.

\subsection{Four Exponentials in disguise}

The \(2\times2\) array
\[
 \begin{pmatrix}
  e^{\log2}&e^{\beta\log2}\\
  e^{\log3}&e^{\beta\log3}
 \end{pmatrix}
 =\begin{pmatrix}2&A\\3&B\end{pmatrix}
\]
is exactly a special Four Exponentials configuration. Any argument whose only transcendence input concerns these four entries, without using additional arithmetic of \(A,B\), is likely to be a reformulation of the open barrier.

\subsection{Adding a ratio without changing the bases}

The reciprocal \(1/\beta\) maps \((A,B)\) back to \((2,3)\), but it is not known to lie in \(\Ealg(2,3)\). Likewise, a fractional-linear expression such as
\[
       \frac{m+\ell\beta}{n+k\beta}
\]
is naturally a common exponent for the changed base pair \(\Pair(n+k\beta)\), not automatically for \((2,3)\). Failing to track the base pair can create a spurious third exponent.

\subsection{A single Hankel identity}

For one triple, the equation
\[
       (\log V)^2=\log U\log W
\]
merely defines a real slope \(\xi=\log V/\log U\). It does not prove \(\xi\in\Q\). Rationality follows in \cref{thm:twohankel} because a second multiplicatively independent initial value supplies the second Six Exponentials row.

\subsection{Separate linear-form bounds}

Baker theory gives strong lower bounds for nonzero expressions such as
\[
       q\log A-p\log2,
       \qquad
       q\log B-p\log3.
\]
For a counterexample these two forms are scalar multiples of the same number \(q\beta-p\). Applying lower bounds separately does not disturb their exact proportionality. A successful use of linear forms would have to create a new, nonproportional form from arithmetic information not already encoded by \(A=2^\beta,B=3^\beta\).

\subsection{Formal \(p\)-adic exponentiation}

The real number \(\beta\) does not automatically define \(2^\beta\) or \(3^\beta\) in a \(p\)-adic field. One may take \(p\)-adic logarithms of suitable rational units built from \(A,B,2,3\), but transferring the real identity to a \(p\)-adic exponential identity requires a separate algebraic relation. The Kummer formulation in \cref{thm:localdefect} is safe because it uses only integer valuations and exact perfect-power classes.

\subsection{Numerical nonvanishing}

For fixed bounds on \(A,B\), interval arithmetic can certify that
\[
       \log A\log3-\log B\log2\ne0
\]
for every tested pair. Such a finite computation is useful for auditing examples but cannot distinguish an open transcendence phenomenon from a very sparse exceptional set. The structural normal form is more useful computationally because it removes every pair with nontrivial defect before expensive logarithmic certification.

\section{New research programs suggested by the normal form}

The exact results narrow the search to a canonical pair satisfying simultaneously
\begin{equation}\label{eq:canonicalresearch}
 \begin{gathered}
  A=2^\beta,\quad B=3^\beta,\quad \beta\ge1\text{ transcendental},\\
  \min(v_2(A),v_3(B))=0,\quad \defect(A,B)=1,\\
  \Erat(2,3)=\Z+\Z\beta,
  \quad
  \Eint(2,3)=\N_0+\N_0\beta.
 \end{gathered}
\end{equation}
The following programs are ordered by how directly they attack a proved bottleneck.

\subsection{Program A: force a Kummer alignment}

For a prime \(\ell\), seek an argument that the real logarithmic identity forces
\[
       [A]=[2]^r,
       \qquad
       [B]=[3]^r
       \quad\text{in }\Q_{>0}^{\times}/\Q_{>0}^{\times\ell}
\]
for some common \(r\). By \cref{thm:localdefect}, any such \(\ell\) contradicts canonical primitivity. This formulation separates the task into two components:
\begin{enumerate}
\item an archimedean-to-Kummer bridge producing congruences of valuation vectors;
\item a local argument showing that at least one prime \(\ell\) must carry the bridge.
\end{enumerate}
The present report supplies the exact endpoint but not the bridge.

\subsection{Program B: nonlinear orbit extension}

Construct an algebraic expression whose exponent is genuinely nonlinear in \(\beta\). The cleanest target is the next iterate
\[
       (A^\beta,B^\beta)=\Pair(\beta^2).
\]
A weaker target is any nonaffine polynomial \(F(\beta)\) for which both powers are algebraic. The fractional-linear groupoid should be treated as the ``trivial closure'': it cannot escape \(\Q+\Q\beta\). A successful construction must use an operation not represented by a ratio of two lattice linear forms.

\subsection{Program C: synchronized primitive divisors}

Use convergents \(p/q\) to create the synchronized differences
\[
       2^p-A^q,
       \qquad
       3^p-B^q.
\]
A useful theorem would couple their primitive prime divisors or force a common residue condition on the valuation vectors of \(A,B\). Separate Zsigmondy-type information about each sequence is not enough; the required result must exploit that both approximation errors are generated by the same \(\varepsilon=q\beta-p\). Formula \eqref{eq:syncasymp} is the archimedean synchronization identity to be combined with such arithmetic information.

\subsection{Program D: prove a closure property that leaks from the monoid}

The exact monoid \(\N_0+\N_0\beta\) is closed under addition and multiplication by nonnegative integers. It is not closed under subtraction, division by integers, or nonlinear polynomial maps. Any theorem showing that the set of common integral exponents enjoys one of these extra closures would force a contradiction. Examples of sufficient leakage include:
\[
       \beta-h\in\Eint(2,3)\quad(h>0),
       \qquad
       \frac{\beta-r}{d}\in\Eint(2,3)\quad(d>1)
\]
when the corresponding coefficients violate canonicality, or \(\beta^2\in\Ealg(2,3)\).

\subsection{Program E: low-support classification in primitive form}

Any bounded-support search should impose \eqref{eq:canonicalresearch} at the outset. In particular:
\begin{itemize}
\item neither \(A\) nor \(B\) is supported only on its axis prime;
\item \(A,B\) are multiplicatively independent;
\item the off-axis valuations and the axis difference have gcd \(1\);
\item at least one of \(A\) and \(B\) is coprime to its axis base.
\end{itemize}
This removes all descendants and repeated perfect-power inflations. A rigorous computation can then enumerate valuation patterns rather than raw integers, use interval logarithms only at the last stage, and report results modulo the unique primitive orbit relation.

\subsection{Program F: seek a functional source of the missing column}

One may search for a function \(F(z)\), built from algebraic data attached to \(2,3,A,B\), with a functional equation forcing algebraicity at both \(z=\beta\) and a nonlinear transform such as \(z=\beta^2\). Mahler-type functional equations, special values of polylogarithms, or algebraic-group logarithms are plausible frameworks, but a valid argument must explicitly show why the value at the nonlinear transform is algebraic. Without that step, the construction remains a restatement of \cref{crit:nonlinear}.

\section{Formalization blueprint}

A Lean development can already formalize most of the new arithmetic structure. The two classical transcendence theorems should be isolated behind clearly named interfaces until suitable library versions are available.

\subsection{Layer 1: elementary real-exponential facts}

Formalize:
\begin{enumerate}
\item positivity and injectivity of \(x\mapsto a^x\) for \(a>1\);
\item rational \(x\) with \(2^x\in\N\) implies \(x\in\Z\);
\item multiplicative independence of \(2\) and \(3\);
\item uniqueness of representations \(n+k\beta\) when \(\beta\notin\Q\).
\end{enumerate}
These should use existing real logarithm, rational power, and unique-factorization material.

\subsection{Layer 2: finite valuation vectors}

Represent a positive rational by a finitely supported map from primes to \(\Z\). For a pair, use a direct sum of two such maps. Define
\[
 g_0=\val(2,3),\qquad g_1=\val(A,B)
\]
and the finite support consisting of \(2,3\) and primes dividing \(AB\). Prove that the gcd of all \(2\times2\) minors is the explicit defect \eqref{eq:defect}. This avoids needing a fully general infinite-dimensional Smith-normal-form API.

\subsection{Layer 3: root extraction and canonicalization}

Prove the equivalence
\[
 d\mid\defect(A,B)
 \quad\Longleftrightarrow\quad
 \exists r,C,D,
 \ A=2^rC^d\ \wedge\ B=3^rD^d.
\]
Define a canonicalization function that:
\begin{enumerate}
\item computes \(d=\defect(A,B)\);
\item chooses the common least residue \(r\);
\item divides valuations by \(d\);
\item subtracts the axis content.
\end{enumerate}
Prove idempotence on canonical pairs and prove the descendant formulas \eqref{eq:descaxis}--\eqref{eq:localizedpower}.

\subsection{Layer 4: abstract transcendence interfaces}

A clean development can use two assumptions with mathematically exact signatures:

\begin{quote}\small
\texttt{six\_exponentials\_rank\_le\_two}: for multiplicatively independent positive algebraic \(u,v\), the \(\Q\)-span of common algebraic exponents has dimension at most two.

\texttt{gelfond\_schneider\_real}: if \(a>0\) is algebraic, \(a\ne1\), and \(x\) is algebraic irrational, then \(a^x\) is transcendental.
\end{quote}

The rest of the report's main theorems should depend on these interfaces rather than on an opaque omnibus axiom.

\subsection{Layer 5: rational and integral exponent classification}

Assuming the rank-two interface, prove
\[
       \Ealg(2,3)=\Q+\Q\beta.
\]
Use the saturation theorem to sharpen this to
\[
       \Erat(2,3)=\Z+\Z\beta,
\]
and then use nonnegativity of valuations plus axis normalization to prove
\[
       \Eint(2,3)=\N_0+\N_0\beta.
\]
The uniqueness of the canonical primitive counterexample is then a short order argument.

\subsection{Layer 6: Hankel and Smith lemmas}

The polynomial identity proof of \cref{thm:globalhankel} can be formalized over \(\Q[X]\): linear polynomials whose product is a square must be associates. The parametrization of
\[
       m_1^2=m_0m_2
\]
by \(d(s^2,rs,r^2)\) is elementary number theory. The determinant-lift Smith form can be proved through determinantal divisors rather than a general reduction algorithm.

\subsection{Suggested theorem dependency graph}

\begin{center}
\begin{tabular}{>{\raggedright\arraybackslash}p{0.27\textwidth} >{\raggedright\arraybackslash}p{0.62\textwidth}}
\toprule
Result & Dependencies\\
\midrule
Exact saturation index & finite valuations, gcd of minors\\
Matched root extraction & saturation index, unique factorization\\
Canonical primitive existence & root extraction, axis shift\\
Algebraic exponent plane & Six Exponentials interface\\
Rational exponent group & algebraic plane, saturated lattice\\
Integral free monoid & rational group, valuation nonnegativity\\
Descendant defect \(=k\) & explicit gcd computation\\
Determinant-lift SNF & determinantal divisors\\
Three-step orbit theorem & Six Exponentials, Gelfond--Schneider\\
Square valuation pattern & elementary divisibility\\
Global Hankel classification & transcendence of \(\beta\), UFD in \(\Q[X]\)\\
Reciprocal duality & free-monoid theorem\\
\bottomrule
\end{tabular}
\end{center}

This separation is important for verification: the new lattice and monoid arguments are fully elementary once the rank-two transcendence ceiling is supplied.

\section{Worked arithmetic models}

The examples in this section illustrate the valuation algebra only. They are \emph{not} asserted to arise from a real exponent satisfying the original problem.

\begin{example}[Full-defect extraction]
Let
\[
 A=2^5 5^6 7^9,
 \qquad
 B=3^2 11^3.
\]
Then
\[
 \defect(A,B)=\gcd(6,9,3,2-5)=3.
\]
The common residue of \(a_2=5\) and \(b_3=2\) modulo \(3\) is \(r=2\), and
\[
 A=2^2(2\cdot5^2\cdot7^3)^3,
 \qquad
 B=3^2(11)^3.
\]
Thus the valuation lattice has saturation index \(3\), and the matched-axis cube root is visible exactly as predicted by \cref{thm:rootextract}.
\end{example}

\begin{example}[A primitive valuation model and a descendant]
Take the formal valuation model
\[
 A=5^2 7^3,
 \qquad
 B=3^4 11^5.
\]
It has axis content \(0\) and
\[
 \defect(A,B)=\gcd(2,3,5,4)=1.
\]
For \(n=4,k=6\), form
\[
 C=2^4A^6,
 \qquad
 D=3^4B^6.
\]
Then \(c(C,D)=4\) and \(\defect(C,D)=6\). The determinant lift is
\[
 \begin{pmatrix}
  1&4&0\\
  0&6&0\\
  0&0&6
 \end{pmatrix},
\]
whose Smith form is \(\operatorname{diag}(1,6,6)\).
\end{example}

\begin{example}[Hankel square pattern]
Let \(d=5\), \(r=3\), and \(s=2\). Then
\[
       d(s^2,rs,r^2)=(20,30,45),
       \qquad
       30^2=20\cdot45.
\]
For any positive rational \(T\), the triple
\[
       (U,V,W)=(T^{20},T^{30},T^{45})
\]
satisfies
\[
       U^3=V^2,
       \qquad
       V^3=W^2,
       \qquad
       (\log V)^2=\log U\log W.
\]
Every prime valuation is a multiple of \((4,6,9)\), scaled here by the corresponding valuation of \(T^5\).
\end{example}

\section{Consolidated theorem: the shape of any counterexample universe}

The main conclusions can be collected into one statement.

\begin{theorem}[Conditional counterexample universe]\label{thm:universe}
If the Alaoglu--Erd\H{o}s conjecture is false, then there exist unique \(\beta\in\R\setminus\Z\) and unique integers
\[
       A=2^\beta,\qquad B=3^\beta
\]
with the following properties.
\begin{enumerate}[label=\textup{(\roman*)}]
\item \(\beta\ge1\) is transcendental, \(A,B\) are multiplicatively independent, and
\[
       \min(v_2(A),v_3(B))=0.
\]
\item The valuation lattice generated by \((2,3)\) and \((A,B)\) is saturated; explicitly,
\[
 \gcd\!\left(
 \{v_p(A):p\ne2\},
 \{v_p(B):p\ne3\},
 v_3(B)-v_2(A)
 \right)=1.
\]
\item The common algebraic exponent space is
\[
       \Ealg(2,3)=\Q+\Q\beta.
\]
\item The common rational exponent group is
\[
       \Erat(2,3)=\Z+\Z\beta.
\]
\item The common integral exponent monoid is
\[
       \Eint(2,3)=\N_0+\N_0\beta.
\]
\item Every nonintegral solution \(\alpha=n+k\beta\) has matched-axis content \(n\), saturation defect \(k\), and localized pair \((A^k,B^k)\).
\item The determinant lift attached to \((1,\alpha)\) relative to \((1,\beta)\) has Smith form \(\operatorname{diag}(1,k,k)\).
\item Every positive rational Hankel triple inside the solution orbit lies on one exponent ray and, after choosing coprime \(r,s\), has the exact valuation pattern \(d(s^2,rs,r^2)\).
\item No irrational orbit ratio can have two successive algebraic iterates on two multiplicatively independent algebraic bases. In particular, algebraicity of \(2^{\beta^2}\) and \(3^{\beta^2}\) would contradict the existence of \(\beta\).
\item For the primitive pair itself,
\[
       \Erat(A,B)=\Z+\Z/\beta,
       \qquad
       \Eint(A,B)=\N_0+\N_0/\beta.
\]
\end{enumerate}
\end{theorem}

\begin{proof}
Combine \cref{lem:elementary,cor:Ealgplane,prop:canonicalexist,thm:rationalgroup,thm:freemonoid,cor:uniqueprimitive,thm:descendant,cor:descSNF,thm:globalhankel,thm:threeorbit,thm:reciprocal}.
\end{proof}

\section{Assessment of progress and the remaining gap}

The report replaces a large and poorly organized search space by a single canonical hypothetical object. It also proves that several apparent ways of generating new examples are only reparametrizations:
\begin{itemize}
\item integer shifts multiply by the known axis pair \((2,3)\);
\item positive lattice multiples raise the primitive pair \((A,B)\) to a power;
\item changing to another orbit base produces only fractional-linear exponents;
\item a rank-three rational Hankel branch necessarily collapses to a rational recurrence with square valuation pattern;
\item all root-extractable descendants are measured by one explicit gcd and return to the same primitive root.
\end{itemize}

The remaining gap can now be stated without ambiguity. One must prove that the canonical conditions in \eqref{eq:canonicalresearch} are inconsistent. The two most direct possibilities are:
\begin{enumerate}
\item force a nontrivial matched Kummer root, contradicting \(\defect(A,B)=1\); or
\item force one genuinely nonlinear common algebraic exponent, contradicting \(\Ealg(2,3)=\Q+\Q\beta\).
\end{enumerate}
The first is arithmetic and valuation-theoretic; the second is transcendence- and orbit-theoretic. The new results show exactly where the two approaches would meet: a rational three-step orbit has the square valuation pattern, while an irrational three-step orbit is forbidden by Six Exponentials and Gelfond--Schneider.

No final contradiction is asserted. The gain is an exact conditional classification, a canonical descent invariant, and several sharply falsifiable proof targets suitable for formalization and further attack.

\appendix

\section{Parametrizing integer Hankel triples}

For completeness, we record the elementary parametrization used above.

\begin{lemma}\label{lem:paramhankel}
Let \(x,y,z\in\N\) satisfy \(y^2=xz\). Then there are unique \(d,r,s\in\N\) with \(\gcd(r,s)=1\) such that
\[
       (x,y,z)=d(s^2,rs,r^2).
\]
The uniqueness is up to simultaneously exchanging \((x,s)\) with \((z,r)\) if the endpoints are not ordered.
\end{lemma}

\begin{proof}
Let \(g=\gcd(x,y)\), write \(x=gu\), \(y=gv\) with \(\gcd(u,v)=1\). The relation gives \(g^2v^2=guz\), hence \(gv^2=uz\). Since \(u\mid v^2\) and \(\gcd(u,v)=1\), one has \(u\mid g\); write \(g=du\). Then \(x=du^2\), \(y=duv\), and the original relation gives \(z=dv^2\). Put \(s=u,r=v\). Coprimality and uniqueness follow from the construction.
\end{proof}

\section{Determinantal-divisor audit}

Let \(M\) be the finite integer matrix whose columns are \(g_0,g_1\). The first column has exactly two nonzero entries, both \(1\). If rows are ordered as
\[
 (\text{first},2),\ (\text{second},3),\
 (\text{first},p\ne2),\ (\text{second},q\ne3),
\]
then its nonzero minors are, up to sign,
\[
       b_3-a_2,\qquad a_p\ (p\ne2),\qquad b_q\ (q\ne3).
\]
This verifies directly that no prime valuation was omitted from \eqref{eq:defect}. In particular, shared primes between \(A\) and \(B\) require no special case: their two valuations occur in separate rows and enter the same gcd list.

For the determinant lift \(\widehat G\), the \(2\times2\) minors are
\[
       \Delta,
       \quad a\Delta,\quad b\Delta,\quad c\Delta,\quad d\Delta,
       \quad 0.
\]
The presence of the minor \(\Delta\) itself is what makes the second determinantal divisor exactly \(|\Delta|\), not a proper multiple.

\section{A compact list of proof obligations for future work}

Any claimed complete proof can be audited against the following questions.
\begin{enumerate}
\item Where is a third \(\Q\)-independent common algebraic exponent obtained, or where is a nontrivial matched root forced?
\item If a new exponent is a ratio, are the bases being changed at the same time?
\item If \(x^2\) or another nonlinear expression is used, where is algebraicity of both corresponding powers proved?
\item If a \(p\)-adic logarithm is invoked, what algebraic identity transfers the real exponent relation into the \(p\)-adic domain?
\item If a perfect power is extracted, why do the base-two and base-three axis residues agree?
\item If numerical separation is used, what theorem makes the search finite?
\item If a rank-three determinant vanishes, is the result more than the rational square-pattern branch classified in \cref{thm:globalhankel}?
\item Is the conclusion incompatible with the canonical free monoid, or does it merely produce another element \(n+k\beta\) already allowed by it?
\end{enumerate}

\begin{thebibliography}{99}

\bibitem{AlaogluErdos1944}
L.~Alaoglu and P.~Erd\H{o}s,
\newblock On highly composite and similar numbers,
\newblock \emph{Transactions of the American Mathematical Society} \textbf{56} (1944), 448--469.

\bibitem{Baker1975}
A.~Baker,
\newblock \emph{Transcendental Number Theory},
\newblock Cambridge University Press, Cambridge, 1975.

\bibitem{BombieriGubler}
E.~Bombieri and W.~Gubler,
\newblock \emph{Heights in Diophantine Geometry},
\newblock Cambridge University Press, Cambridge, 2006.

\bibitem{Cassels}
J.~W.~S.~Cassels,
\newblock \emph{An Introduction to Diophantine Approximation},
\newblock Cambridge University Press, Cambridge, 1957.

\bibitem{Lang1966}
S.~Lang,
\newblock \emph{Introduction to Transcendental Numbers},
\newblock Addison--Wesley, Reading, Massachusetts, 1966.

\bibitem{Newman}
M.~Newman,
\newblock \emph{Integral Matrices},
\newblock Academic Press, New York, 1972.

\bibitem{NivenZuckermanMontgomery}
I.~Niven, H.~S.~Zuckerman, and H.~L.~Montgomery,
\newblock \emph{An Introduction to the Theory of Numbers}, 5th ed.,
\newblock Wiley, New York, 1991.

\bibitem{Waldschmidt2000}
M.~Waldschmidt,
\newblock \emph{Diophantine Approximation on Linear Algebraic Groups: Transcendence Properties of the Exponential Function in Several Variables},
\newblock Grundlehren der mathematischen Wissenschaften, vol.~326, Springer, Berlin, 2000.

\bibitem{WaldschmidtSurvey}
M.~Waldschmidt,
\newblock Open Diophantine problems,
\newblock \emph{Moscow Mathematical Journal} \textbf{4} (2004), 245--305.

\end{thebibliography}

\end{document}
